\section{Optimization Details}
\label{appendix:optimization_hps}

\subsection{Notation}
We consider the standard multi-agent reinforcement learning formalism of $N$ agents interacting with each other in an environment. This interaction is defined by a set of states $\mathcal{S}$ describing the state of the world and configurations of all agents, a set of observations $\mathcal{O}^1, \ldots \mathcal{O}^N$ of all agents, a set of actions $\mathcal{A}^1, \ldots, \mathcal{A}^N$ of all agents, a transition function $\mathcal{T}: \mathcal{S} \times \mathcal{A}^1 \ldots \mathcal{A}^N \rightarrow \mathcal{S}$ determining the distribution over next states, and a reward for each agent $i$ which is a function of the state and the agent's action.
Agents choose their actions according to a stochastic policy $\pi_{\theta_i}: \mathcal{O}^i \times \mathcal{A}^i \rightarrow [0, 1]$, where $\theta_i$ are the parameters of the policy. %and $\mathcal{S}^I$ is a set of internal states, such as the hidden state of a recurrent network.
In our formulation, policies are shared between agents, $\pi_{\theta_i} = \pi_{\theta}$ and the set of observations $\mathcal{O}$ contains information for which role (e.g. hider or seeker) the agent will be rewarded. Each agent $i$ aims to maximize its total expected discounted return $R^i = \sum_{t=0}^H \gamma^t r_t^i$, where $H$ is the horizon length and $\gamma$ is a time discounting factor that biases agents towards preferring short term rewards to long term rewards. The \textit{action-value function} is defined as $Q^{\pi_i}(s_t,a_t^i) = \mathbb{E}[R^i_t|s_t,a_t^i]$, while the \textit{state-value function} is defined as $V^{\pi_i}(s_t) = \mathbb{E}[R_t|s_t]$. The \textit{advantage function} $A^{\pi_i}(s_t,a_t^i) := Q^{\pi_i}(s_t,a_t^i) - V^{\pi_i}(s_t)$ describes whether taking action $a_t^i$ is better or worse for agent $i$ when in state $s_t$ than the average action of policy $\pi_i$.

\subsection{Proximal Policy Optimization (PPO)}
Policy gradient methods aim to estimate the gradient of the policy parameters with respect to the discounted sum of rewards, which is often non-differentiable.
A typical estimator of the policy gradient is $g := \mathbb{E}[\hat{A}_t \nabla_\theta \log \pi_\theta]$, where $\hat{A}$ is an estimate of the advantage function. PPO \citep{schulman2017proximal}, a policy gradient variant, penalizes large changes to the policy to prevent training instabilities. PPO optimizes the objective $L = \mathbb{E}\left[\min(l_t(\theta)\hat{A}_t, \textrm{clip}(l_t(\theta), 1 - \epsilon, 1 + \epsilon)\hat{A}_t\right]$, where $l_t(\theta) = \frac{\pi_\theta(a_t|s_t)}{\pi_{old}(a_t|s_t)}$ denotes the likelihood ratio between new and old policies and $\textrm{clip}(l_t(\theta), 1 - \epsilon, 1 + \epsilon)$ clips $l_t(\theta)$ in the interval $[1-\epsilon, 1+\epsilon]$.

\subsection{Generalized advantage estimation}
We use Generalized Advantage Estimation \citep{schulman2015high} with horizon length $H$ to estimate the advantage function. This estimator is given by:
\begin{equation*}
    \hat{A}_t^H = \sum_{l=0}^H (\gamma \lambda)^l \delta_{t+l}, \qquad \delta_{t+l} :=  r_{t+l} + \gamma V(s_{t+l+1}) - V(s_{t+l}) 
\end{equation*}
where $\delta_{t+l}$ is the TD residual, $\gamma, \lambda \in [0,1]$ are discount factors that control the bias-variance tradeoff of the estimator, $V(s_t)$ is the value function predicted by the value function network and we set $V(s_t) = 0$ if $s_t$ is the last step of an episode. This estimator obeys the reverse recurrence relation $\hat{A}_t^H = \delta_t + \gamma \lambda \hat{A}_{t+1}^{H-1}$.

We calculate advantage targets by concatenating episodes from policy rollouts and truncating them to windows of $T=160$ time steps (episodes contain 240 time steps). If a window $(s_0, \ldots, s_{T-1})$ was generated in a single episode we use the advantage targets $(\hat{A}_0^{H=T}, \hat{A}_1^{H=T-1}, \ldots, \hat{A}_{T-1}^{H=1})$. If a new episode starts at time step $j$ we use the advantage targets $(\hat{A}_0^{H=j}, \hat{A}_1^{H=j-1}, \ldots, \hat{A}_{j-1}^{H=1}, \hat{A}_j^{H=T-j}, \ldots, \hat{A}_{T-1}^{H=1})$.

Similarly we use as targets for the value function $(\hat{G}_0^{H=T}, \hat{G}_1^{H=T-1}, \ldots, \hat{G}_{T-1}^{H=1})$ for a window generated by a single episode and $(\hat{G}_0^{H=j}, \hat{G}_1^{H=j-1}, \ldots, \hat{G}_{j-1}^{H=1}, \hat{G}_j^{H=T-j}, \ldots, \hat{G}_{T-1}^{H=1})$ if a new episode starts at time step $j$, where the return estimator is given by $\hat{G}_t^H := \hat{A}_t^H + V(s_t)$. This value function estimator corresponds to the TD($\lambda$) estimator \citep{sutton2018}.

\subsection{Normalization of observations, advantage targets and value function targets}
We normalize observations, advantage targets and value function targets. Advantage targets are z-scored over each buffer before each optimization step. Observations and value function targets are z-scored using a mean and variance estimator that is obtained from a running estimator with decay parameter $1 - 10^{-5}$ per optimization substep.


\subsection{Optimization setup}
Training is performed using the distributed \textit{rapid} framework \citep{openai2018}. Using current policy and value function parameters, CPU machines roll out the policy in the environment, collect rewards, and compute advantage and value function targets. Rollouts are cut into windows of 160 timesteps and reformatted into 16 chunks of 10 timesteps (the BPTT truncation length). The rollouts are then collected in a training buffer of 320,000 chunks. Each optimization step consists of 60 SGD substeps using Adam with mini-batch size 64,000. One rollout chunk is used for at most 4 optimization steps. This ensures that the training buffer stays sufficiently on-policy.

\subsection{Optimization hyperparameters}\label{appendix:hyperparam}

Our optimization hyperparameter settings are as follows:

\begin{center}
\begin{tabular}[h]{|c|c|}
    \hline
	Buffer size & 320,000 \\
	\hline
	Mini-batch size & 64,000 chunks of 10 timesteps\\
	\hline
	Learning rate & $3 \cdot 10^{-4}$ \\
	\hline
	PPO clipping parameter $\epsilon$ & 0.2 \\
	\hline
	Gradient clipping & 5 \\
	\hline
	Entropy coefficient &  0.01 \\
	\hline
	$\gamma$ & 0.998  \\
	\hline
	$\lambda$ & 0.95 \\
	\hline
	Max GAE horizon length $T$ & 160 \\
	\hline
	BPTT truncation length & 10 \\
	\hline
\end{tabular}
\end{center}

\subsection{Policy architecture details}

Lidar observations are first passed through a circular 1D-convolution and concatenated onto the agents representation of self, $x_\text{self}$. Each object is concatenated with $x_\text{self}$ and then embedded with a dense layer where parameters are shared between objects of the same type, e.g. all boxes share the same embedding weights.
All the embedded entities are then passed through a residual self-attention block, similar to \cite{vaswani2017attention} but without position embeddings, in the form of $y = \text{dense}(\text{self\_attention}(x)) + x$.
We then average-pool entity embeddings and concatenate this pooled representation to $x_\text{self}$. Note that in the policy network the entities not observed by each agent are masked away through self-attention and pooling. Finally, this pooled representation is passed through another dense layer and an LSTM \citep{hochreiter1997long} before pulling off separate action heads for each of the 3 action types described in Section~\ref{sec:environment}. 
We also add layer normalization~\citep{ba2016layer} to every hidden layer of the policy network except the 1D-convolution layer. We empirically observe that layer normalization leads to faster training and better transfer performance.

\begin{center}
\begin{tabular}[h]{|c|c|}
    \hline
	Size of embedding layer & 128 \\
	\hline
	Size of MLP layer& 256 \\
	\hline
	Size of LSTM layer & 256 \\
	\hline
	Residual attention layer & 4 attention heads of size 32 \\
	\hline
	Weight decay coefficient & $10^{-6}$ \\
	\hline
\end{tabular}
\end{center}
